\subsection{Limitations}
\label{sec:limitations}

\paragraph{One model family and one harness.} All subjects, judges, and reviewers are Claude models, and all agents ran in the Claude Code harness. We chose this to compare generations within a family under a production scaffold, but it means we cannot say how models from other developers report their own gaming. Judges and reviewers from the same family as the subjects may share blind spots; the behavioral outcomes (tampering, special-casing, hidden-test failures) are computed without any model, and we validated the report coding against blinded human labels, but the review experiment has no human reviewers.

\paragraph{Small, constructed repositories.} Each task has a single conflicting test in a short function, and several docstrings state the rule the test violates outright (for example, ``the cap always wins''), which makes the conflict unambiguous but also more conspicuous than in real code. Real repositories contain more context, more ambiguity, and more ways to be wrong. We expect the pressure toward rationalization to grow, not shrink, with ambiguity, but we did not test this.

\paragraph{We cannot see the reasoning.} The harness returns no reasoning text, so we cannot tell whether an agent believed its invented rule when it wrote it. The follow-up probes show what the agent says when asked, which bounds but does not settle the question of belief. We therefore describe the reports as fabricated with respect to the repository, not as intentionally deceptive.

\paragraph{The minimal pairs change two things at once.} Replacing an explainable expected value with one that no plausible rule produces also makes the test error easier to recognize, and in some variants the new value is visibly impossible (for example, allocating seven of six desks, or a misspelled word). We regard these as two descriptions of the same property from the agent's point of view, since a value that no rule can produce is exactly one that reads as a mistake, but the experiment does not separate them.

\paragraph{Sample sizes.} We ran two or three trials per model, task, and condition. Pooled estimates are reasonably tight, but per-model intervals are wide, and our budget did not allow more trials for the most expensive models. Trials that timed out (15 minutes) were excluded and counted; long searches for an explanatory rule are one plausible cause, so exclusions may slightly understate gaming.

\paragraph{Hidden-test coverage.} Two tasks call the conflict input with keyword arguments in the hidden tests, so a special case keyed to positional arguments would not register as collateral damage there. This affects the collateral-damage measure only, not the detection of special-casing.
